\chapter{Machine Learning Based Reconstruction}
\label{chap:ml-reconstruction}


In addition to the outline reconstruction methods described in \Cref{chap:reconstruction-methods}, this chapter focuses on the development of a machine learning-based reconstruction strategies for the dileptonic \ttbar channel. Machine learning-based reconstruction techniques have gained significant traction in recent years. One major example for this, is their use in flavour tagging algorithms \cite{flavour_tagging_review}. For this, machine learning models are trained to address both challenges of jet assignment and neutrino regression, by learning from simulated data to predict the correct jet-parton assignments and neutrino momenta.

For the high permutation complexity of jet assignments in all-hadronic channels, SPANet \cite{spanet} has demonstrated significant improvements over traditional methods. These architectures leverage attention mechanisms to improve the accuracy by leveraging the full event information and utilising permutation equivariant architectures to handle the combinatorial nature of the problem.
Based on these successes, this thesis explores the adaptation of similar architectures for the dileptonic \ttbar channel, which presents unique challenges due to the presence of two neutrinos and fewer jets.

On the other hand, for neutrino regression, $\nu^2$-Flows \cite{nu_square_flows} has shown promise in improving the reconstruction of neutrino momenta by modelling the underlying probability distributions of the neutrino kinematics.

This thesis investigates, how machine learning can be used to provide a full reconstruction pipeline for the dileptonic \ttbar channel, by combining jet assignment and neutrino regression.

\section{Jet-Assignment}
\label{sec:jet-assignment}

\subsection{Neural Network Architectures}
\label{sec:nn-architectures}


\subsection{Training Procedure}
\label{sec:training-procedure}

\subsection{Hyperparameter Optimisation}
\label{sec:hyperparameter-optimisation}

\section{Neutrino Regression}
\label{sec:neutrino-regression}

